\section{相关工作}

传统 RL 框架如 verl~\cite{hybridflow}、AReaL~\cite{areal} 与 slime~\cite{slime} 最初要求智能体循环直接实现在训练框架内部，遵循经典 ReAct 式~\cite{react} 马尔可夫形式——训练引擎拥有环境交互循环。这使得复用独立维护的智能体脚手架（如 mini-SWE-agent~\cite{minisweagent}、OpenHands~\cite{openhands}、OpenCode~\cite{opencode}、Claude Code~\cite{claudecode}、Codex~\cite{codex}、OpenClaw~\cite{openclaw} 与 Hermes~\cite{hermes}）十分困难，因为每个都需要在训练栈中重新实现。我们最初的 Agent Lightning~\cite{agentlightning} 工作提出了解耦架构，转而通过 LLM 端点将任意智能体脚手架接入 RL 训练；这种基于代理的方式此后被 verl Uni-Agent~\cite{uniagent}、AReaL 2.0~\cite{areal2}、slime v0.3.0~\cite{slime} 与 Polar~\cite{polar} 采纳。如第~\ref{sec:challenges} 节所述，这些框架在重分词、优势计算与损失归一化（每次 rollout 产生动态数量训练样本）上做出不同甚至相互冲突的设计选择，并普遍依赖商业沙箱服务（如 Modal Sandbox、Volcano veFaas 与 E2B）大规模执行智能体。Agent Lightning v1.0 则完全运行在自托管 Kubernetes 集群上，整个系统约 3500 行代码实现，为研究这些设计选择提供了紧凑透明的试验平台。我们分别在 Search-R1~\cite{searchr1}、LLM-in-Sandbox~\cite{llminsandbox} 与 SWE-smith~\cite{swesmith} 的搜索智能体、指令遵循与编程智能体设置上验证了它。
